\chapter{Conclusion and Future Work}
\label{chap:conclusion}

%% -----------------------------------------------------------------------
\section{Summary of Contributions}
\label{sec:conc:summary}

This DDP has developed, validated, and deployed a digital-twin simulation
framework for Multi-Echelon Inventory Optimisation (MEIO) with
transport-aware joint optimisation.  The main contributions are:

\begin{enumerate}
  \item \textbf{A modular, open digital-twin framework} implementing
        seven inventory control policies (base-stock, $(s,S)$, $(R,Q)$,
        $(R,S)$, order-up-to with phase offset, $(k,m)$-cycle, and
        echelon base-stock), an arbitrary DAG-topology network model,
        and a transport layer with simultaneous volume and weight capacity
        constraints.  The framework is validated by approximately 80
        pytest tests covering simulator invariants, policy behaviour, and
        transport mechanics.

  \item \textbf{An Optuna-TPE joint optimiser} that simultaneously
        searches over 15 integer base-stock levels and 5 continuous
        dispatch thresholds in a 20-dimensional mixed-integer space,
        subject to a configurable minimum fill-rate constraint.

  \item \textbf{Empirical demonstration of the 39\% cost-reduction
        achievable through joint inventory-transport optimisation.}  On
        a 5-SKU, 1N3 network with M5 Walmart demand, the joint optimiser
        (E3) achieves:
        \begin{itemize}
          \item Total cost \$1,556,516 vs \$2,553,418 at baseline
                ($-$39.0\%, saving nearly \$1M).
          \item Fill rate 95.37\% vs 86.15\% at baseline (+9.2pp),
                comfortably above the 92\% service-level target.
          \item Transport cost reduced from 74.1\% to 54.4\% of total,
                driven by a 4.3:1 leverage ratio between holding and
                transport savings.
        \end{itemize}

  \item \textbf{A Pareto frontier mapping the cost-service trade-off}
        across eight fill-rate targets (80\%--98\%), revealing a natural
        operating point at $\sim$95\% fill / \$1.56M and sharply
        diminishing returns above 94\%.

  \item \textbf{A CSV-driven configuration pipeline} lowering the barrier
        for practitioners to define and experiment with new scenarios
        without editing JSON directly.

  \item \textbf{Preliminary validation experiments (EV1--EV4)} confirming
        simulator correctness against known theoretical behaviours: order
        convergence under constant demand, demand shock propagation, policy
        comparison under transport capacity, and non-linear capacity
        scaling.
\end{enumerate}

%% -----------------------------------------------------------------------
\section{Future Work}
\label{sec:conc:futurework}

The framework is designed to be extensible.  Several directions are
planned or proposed for future development.

\begin{description}
  \item[Supply disruption sensitivity (near-term).] The supply disruption
        experiment (§7.9) will re-run E0 and E3 under a 14-day supplier
        outage to quantify how the optimised configuration degrades relative
        to the analytical baseline.  The hypothesis is that E3's elevated
        warehouse inventory provides a meaningful buffer.

  \item[Stochastic lead times.] The current simulator uses deterministic
        transit times.  Extending to stochastic lead times (e.g.\ Gamma
        or Uniform distributed) would allow the framework to model
        carrier variability and quantify its impact on optimal safety stock.

  \item[Demand forecasting integration.] Replacing the CSV demand reader
        with a demand forecasting module (e.g.\ Exponential Smoothing,
        LightGBM, or neural-network-based point and interval forecasts)
        would enable evaluation of policies under realistic forecasting
        error rather than perfect demand information.

  \item[Lateral transshipment.] Allowing stock transfers between
        co-located retailers can substantially reduce local stockouts.
        Adding a transshipment option to the edge model and a
        transshipment-aware policy would extend the framework to cover
        this common industry practice.

  \item[RL-based policies.] Reinforcement learning (RL) policies
        (e.g.\ Proximal Policy Optimisation) could be trained within the
        simulator to discover adaptive ordering strategies that the
        analytical policy zoo cannot represent.

  \item[Multi-objective Pareto optimisation.] Rather than penalising
        fill-rate violations as a scalar, a future version could use
        Optuna's multi-objective API to jointly minimise cost and
        maximise fill rate, yielding a richer Pareto front with
        continuous trade-off resolution.

  \item[Web dashboard.] A lightweight web interface allowing practitioners
        to configure scenarios, trigger optimisation runs, and explore
        results without command-line interaction would broaden the
        framework's accessibility.

  \item[Industry-specific configuration templates.] Pre-built JSON and
        CSV templates for common supply-chain archetypes (grocery retail,
        pharmaceutical distribution, e-commerce fulfilment) would reduce
        setup time for new users.
\end{description}

%% -----------------------------------------------------------------------
\section{Closing Remarks}
\label{sec:conc:closing}

The central finding of this DDP is simple to state but non-obvious in
practice: \emph{inventory and transport are not independent levers.}
Holding extra safety stock is conventionally viewed as a cost to be
minimised; this work shows that, in a transport-cost-dominated network,
extra holding cost is a strategic investment that enables vehicle
consolidation and generates a 4:1 return in transport savings.

The digital-twin simulation framework built in this project provides
a portable, testable, and configurable tool for discovering such
counter-intuitive optima in real supply chains.  By exposing the
full cost-service Pareto frontier, it enables data-driven negotiation
of service-level agreements and moves supply-chain design from intuition
to evidence.
